\section{Data}
\label{sec:data}

% ROTEIRO: ANTAQ (Estatistico Aquaviario) e Comex Stat; construcao do painel
% porto--mes; conceitos oficiais de tempo; crosswalk URF--porto e ressalvas.

The operational panel combines \Natracacoes\ vessel calls with cargo
movement recorded by ANTAQ between \Anoinicio\ and \Anofim, aggregated into
\Npainellinhas\ port--month observations covering \Nportos\ ports
(\Nportospublicos\ public ports and \Nportostups\ private terminals).

\input{tables/tab01_amostra}

Operational times follow the official definitions of the agency: $T_1$ is the
interval between arrival at the anchorage and berthing; $T_2$ the wait between
berthing and the start of operations; $T_3$ the cargo operation itself; and
$T_4$ the wait until unberthing, with $T_A = T_2+T_3+T_4$ and $T_E = T_1+T_A$.

\input{tables/tab02_tempos}

\begin{figure}[htbp]
\centering
\includegraphics[width=\textwidth]{fig01_distribuicao_t1.pdf}
\caption{Components of the vessel's stay in port, by year and port type.
Median hours per vessel call with cargo movement (deep-sea, cabotage and
inland navigation). The shaded band marks the single-window rollout
(July 2011--April 2013). Source: ANTAQ waterway statistics, own elaboration.}
\label{fig:t1dist}
\end{figure}

\begin{figure}[htbp]
\centering
\includegraphics[width=\textwidth]{fig11_map.pdf}
\caption{Port facilities in the panel by treatment status. Marker area is proportional to
tonnes handled in 2010--2015. Source: ANTAQ waterway statistics, own elaboration. State boundaries: IBGE (2024).}
\label{fig:map}
\end{figure}

\input{tables/tab04_qualidade}

\FloatBarrier
